%% ============================================================
%% AULA 7 — Redes que enxergam: convolução
%% GBC073 — Inteligência Computacional (FACOM/UFU)
%%
%% Esqueleto com esboço de conteúdo — a desenvolver.
%% Ementa (ficha SEI 5116656): item 2 — arquiteturas;
%% aplicações: classificação.
%%
%% Compilar:  latexmk -xelatex aula07.tex   (duas passadas!)
%% ============================================================
\documentclass[aspectratio=169, 11pt]{beamer}

\makeatletter
\def\input@path{{./}{../}}
\makeatother

\usetheme{InteligenciaComputacional}   % ANTES do babel: fontspec vem primeiro
\usepackage{babel}
\babelprovide[import, main]{brazilian}

\title[Aula 7 — Convolução]{Aula 7: Redes que enxergam — convolução}
\subtitle[GBC073 — Inteligência Computacional]{Campos receptivos \textbullet\ pooling \textbullet\ LeNet a ResNet \textbullet\ transferência}
\author[Prof.\ Marcelo Albertini]{Prof.\ Marcelo Keese Albertini}
\institute{GBC073 — Inteligência Computacional \textbullet\ Universidade Federal de Uberlândia}
\date{2026}

\begin{document}

% ------------------------------------------------------------ capa
\begin{frame}[plain, noframenumbering]
  \titlepage
\end{frame}

% ------------------------------------------------------------ roteiro
\begin{frame}{Roteiro da aula}
  \begin{itemize}\setlength{\itemsep}{0.5ex}
    \item<1-> \textbf{O que a convolução muda em relação à MLP}
    \item<2-> \textbf{Filtros e mapas de características}
    \item<3-> \textbf{Pooling e invariância}
    \item<4-> \textbf{Arquiteturas clássicas e transferência de aprendizado}
  \end{itemize}
\end{frame}

% ------------------------------------------------------------
\section{O que a convolução muda em relação à MLP}

\begin{frame}{O que a convolução muda em relação à MLP}
  \begin{itemize}\setlength{\itemsep}{0.4ex}
    \item Uma MLP sobre pixels ignora a estrutura espacial da imagem.
    \item Convolução: o mesmo filtro desliza por toda a imagem — pesos
          compartilhados.
    \item Duas hipóteses embutidas: localidade (campos receptivos) e
          invariância à translação.
    \item Muito menos parâmetros, muito mais dados aproveitados.
  \end{itemize}
\end{frame}

% ------------------------------------------------------------
\section{Filtros e mapas de características}

\begin{frame}{Filtros e mapas de características}
  \begin{itemize}\setlength{\itemsep}{0.4ex}
    \item Um filtro aprende um padrão local: bordas, cantos, texturas.
    \item Camadas profundas compõem padrões em padrões.
    \item \emph{Stride}, \emph{padding}, canais — a aritmética da convolução.
    \item Visualizar os filtros aprendidos: o que a rede vê no mundo?
  \end{itemize}
\end{frame}

% ------------------------------------------------------------
\section{Pooling e invariância}

\begin{frame}{Pooling e invariância}
  \begin{itemize}\setlength{\itemsep}{0.4ex}
    \item Pooling resume uma região: máximo ou média.
    \item Reduz a resolução espacial e ganha tolerância a pequenos deslocamentos.
    \item Efeito colateral: a rede fica mais barata a cada camada.
    \item Invariância total não existe — é sempre aproximada e aprendida.
  \end{itemize}
\end{frame}

% ------------------------------------------------------------
\section{Arquiteturas clássicas e transferência de aprendizado}

\begin{frame}{Arquiteturas clássicas e transferência de aprendizado}
  \begin{itemize}\setlength{\itemsep}{0.4ex}
    \item Linha do tempo: LeNet, AlexNet, VGG, Inception, ResNet.
    \item ResNet: conexões residuais que permitem redes muito profundas.
    \item Transferência: reutilizar redes pré-treinadas e ajustar
          (\emph{fine-tuning}).
    \item Poucos dados + grande modelo pré-treinado = combinação vencedora.
  \end{itemize}
  \begin{alertabox}
    Treinar uma CNN do zero em um conjunto pequeno quase sempre perde para
    o \emph{fine-tuning} de uma rede pré-treinada — comece por lá.
  \end{alertabox}
\end{frame}

\end{document}
